\input{preamble}

\title{Section 6.1 --- Coordinate Geometry --- Answer Key}

\begin{document}
\maketitle

\section*{Practice With Points}
\begin{senum}
\item On a separate sheet of graph paper (or the other side of this sheet), draw appropriate axes and:
\begin{nmenum}
\mitemxx
{\begin{GridFilled}[0.15in]{4}{2}{5}{3}
\foreach \x/\y/\lbl in {-4/3/A,2/-5/B} {
	\node[inner sep=0pt] (\lbl) at (\x,\y){};
	\filldraw[black] (\lbl) circle (\dotsize);
	\draw (\lbl) node[anchor=south west]{$\lbl$};
}
\draw[thick] (A) -- (B);
\end{GridFilled}}
{There are two possible diagrams you could use here.\\
\begin{ilpic}[x=0.15in,y=0.15in]
\begin{scope}[xshift=-4*0.15in]
\DrawGrid{4}{2}{5}{3}
\foreach \x/\y/\lbl in {-4/3/A,2/-5/B,-4/-5/C} {
	\node[inner sep=0pt] (\lbl) at (\x,\y){};
	\filldraw[black] (\lbl) circle (\dotsize);
	\draw (\lbl) node[anchor=south west]{$\lbl$};
}
\draw[thick] (A) -- (B);
\draw[thick,dashed] (A) -- (C);
\draw[thick,dashed] (B) -- (C);
\draw (-4.9,-7.4) node[anchor=south west]{Diagram I};
\end{scope}
\begin{scope}[xshift=5*0.15in]
\DrawGrid{4}{2}{5}{3}
\foreach \x/\y/\lbl in {-4/3/A,2/-5/B,2/3/C} {
	\node[inner sep=0pt] (\lbl) at (\x,\y){};
	\filldraw[black] (\lbl) circle (\dotsize);
	\draw (\lbl) node[anchor=south west]{$\lbl$};
}
\draw[thick] (A) -- (B);
\draw[thick,dashed] (A) -- (C);
\draw[thick,dashed] (B) -- (C);
\draw (-4.9,-7.4) node[anchor=south west]{Diagram II};
\end{scope}
\end{ilpic}
}
\mitemx
{With Diagram I: $C=(-4,-5)$\\With Diagram II: $C=(2,3)$}
\mitemx
{With Diagram I: $AC=8$\\With Diagram II: $AC=6$}
\mitemx
{With Diagram I: $BC=6$\\With Diagram II: $BC=8$}
\mitemxx
{$AB=10$}
{$M=(-1,1)$%
\begin{ilpic}[x=0.15in,y=0.15in]
\useasboundingbox(-5,-6) rectangle (3,-2);
\DrawGrid{4}{2}{5}{3}
\foreach \x/\y/\lbl/\dir in {-4/3/A/{south west},2/-5/B/{south west},-1/-1/M/{north east}} {
	\node[inner sep=0pt] (\lbl) at (\x,\y){};
	\filldraw[black] (\lbl) circle (\dotsize);
	\draw (\lbl) node[anchor=\dir]{$\lbl$};
}
\draw[thick] (A) -- (B);
\end{ilpic}}
\end{nmenum}
\end{senum}

\section*{Distance Formula}
Using the distance formula (without plotting on a graph) find the distance between the two points given. If necessary, give your answer as a simplified radical and as a decimal approximation.
\begin{smenum}
\mitemxxx
{$13$}
{$5$}
{$\sqrt{5}\approx 2.236$}
\mitemxxx
{$2.9$}
{$1$}
{$\frac{2\sqrt{298}}{9}\approx 3.836$}
\end{smenum}

\section*{Midpoint Formula}
For each pair of points $A$ and $B$ given below, find the midpoint of $\overline{AB}$.
\begin{smenum}
\mitemxxx
{$M=(2,2)$}
{$M=(8,4.5)$}
{$M=(-8.5,-4)$}
\mitemxxx
{$M=\left(\frac{5}{2},3\right)$}
{$M=(1.3,3.7)$}
{$M=(-38.5,60.5)$}
\end{smenum}

\section*{Triangles and Inequalities}
For each list of three values, classify them as one of the following.
\begin{Clist}
\item Not possible as the distances between three points.
\item Distances between three collinear points.
\item Side lengths of an obtuse triangle.
\item Side lengths of a right triangle.
\item Side lengths of an acute triangle.
\end{Clist}
\begin{smenum}
\mitemxxxxx
{C}
{B}
{E}
{A}
{D}
\mitemxxxxx
{B}
{E}
{C}
{A}
{D}
\end{smenum}

\section*{Classifying Triangles}
Use the distance formula to classify each set of three points as:
\begin{Clist}
\item Three collinear points
\item Corners of an obtuse triangle
\item Corners of a right triangle
\item Corners of an acute triangle
\end{Clist}
\begin{smenum}
\mitemxxxx
{A}
{C}
{B}
{B}
\end{smenum}
\end{document}